\providecommand\ExamFontPath{../../../字体/}
\documentclass[fontset=exam]{exam-zh}
\usepackage{longtable,booktabs,array,calc,multirow,fvextra}
\usepackage[normalem]{ulem}
\providecommand{\hl}[1]{#1}
\providecommand{\ul}[1]{\uline{#1}}
\DefineVerbatimEnvironment{verbatim}{Verbatim}{breaklines=true,fontsize=\small}
\providecommand{\tightlist}{\setlength{\itemsep}{0pt}\setlength{\parskip}{0pt}}
\newcommand{\pandocbounded}[1]{#1}
\setlength{\emergencystretch}{3em}
\setlength{\tabcolsep}{2pt}
\examsetup{page/size=a4paper}
\title{厦门大学 2013 级期中考试 A 卷}
\subject[4em]{计算方法}
\begin{document}
\maketitle
\fancyfoot[C]{第 \thepage 页（共 \pageref*{LastPage} 页）}
\begin{enumerate}
\def\labelenumi{\arabic{enumi}.}
\item
  已知下列数值表，求符合表值的插值多项式，并给出插值余项的表达式。
\end{enumerate}

\begin{center}\begin{tabular}{c|ccc}$x_i$&0&1&2\\$y_i$&2&1&2\\$y_i^{\prime}$&$-2$&$-1$&\\$y_i^{\prime\prime}$&$-10$&&\end{tabular}\end{center}

2. 设\(x_{0},x_{1},\cdots,x_{n}\)为\(n + 1\)个互异的插值节点, \(l_{i}(x),(i = 0,1,\cdots,n)\)为拉格朗日插值基函数.证明:

\[{\sum\limits_{i = 0}^{n}{l_{i}(0)x_{i}^{k}}} = \left\{ \begin{matrix}
1 & {k = 0} \\
0 & {k = 1,2,\cdots,n} \\
{{( - 1)}^{n}x_{0}x_{1}\cdots x_{n}} & {k = n + 1}
\end{matrix} \right.\]

\begin{enumerate}
\def\labelenumi{\arabic{enumi}.}
\setcounter{enumi}{2}
\tightlist
\item
\end{enumerate}

\begin{longtable}[]{@{}
  >{\raggedright\arraybackslash}p{(\linewidth - 8\tabcolsep) * \real{0.1980}}
  >{\raggedright\arraybackslash}p{(\linewidth - 8\tabcolsep) * \real{0.1993}}
  >{\raggedright\arraybackslash}p{(\linewidth - 8\tabcolsep) * \real{0.2008}}
  >{\raggedright\arraybackslash}p{(\linewidth - 8\tabcolsep) * \real{0.2009}}
  >{\raggedright\arraybackslash}p{(\linewidth - 8\tabcolsep) * \real{0.2009}}@{}}
\toprule\noalign{}
\begin{minipage}[b]{\linewidth}\raggedright
x\textsubscript{i}
\end{minipage} & \begin{minipage}[b]{\linewidth}\raggedright
1
\end{minipage} & \begin{minipage}[b]{\linewidth}\raggedright
2
\end{minipage} & \begin{minipage}[b]{\linewidth}\raggedright
3
\end{minipage} & \begin{minipage}[b]{\linewidth}\raggedright
4
\end{minipage} \\
\midrule\noalign{}
\endhead
\bottomrule\noalign{}
\endlastfoot
y\textsubscript{i} & 4 & 10 & 18 & 26 \\
\end{longtable}

用\(y = c_{0} + c_{1}x + c_{2}x^{2}\)拟合上述数据，采用正交多项式方法求解。

\begin{enumerate}
\def\labelenumi{\arabic{enumi}.}
\setcounter{enumi}{3}
\item
  求\(f(x) = x^{4} + 3x^{3} - 1\)在区间{[}0,1{]}上的三次最佳一致逼近多项式。
\item
  确定求积公式
\end{enumerate}

\[{\int_{x_{0}}^{x_{1}}{(x - x_{0})f(x)dx}} = h^{2}\lbrack Af(x_{0}) + Bf(x_{1})\rbrack + h^{3}\lbrack Cf'(x_{0}) + Df'(x_{1})\rbrack + R(f)\]

中的系数\emph{A, B, C, D}, 使其代数精度尽量高，并给出\emph{R}(\emph{f})的表达式。其中\(h = x_{1} - x_{0}\)。

6. \(\left\{ \phi_{n}(x) \right\}\)是{[}a, b{]}上带权ρ(x)的正交多项式序列，x\textsubscript{i} (i=0, 1, \ldots, n)是\(\phi_{n + 1}(x)\)的零点，\(l_{i}(x)(i = 0,1,...,n)\)是以\(\left\{ x_{i} \right\}\)为节点的拉格朗日插值基函数。

\({\int_{a}^{b}{\rho(x)f(x)dx}} \approx {\sum\limits_{k = 0}^{n}{A_{k}f(x_{k})}}\)为高斯求积公式。证明：

\[{\sum\limits_{k = 0}^{n}{\int_{a}^{b}{\rho(x)l_{k}^{2}(x)dx}}} = {\int_{a}^{b}{\rho(x)dx}}.\]

7．用龙贝格方法计算\(\int_{0.3}^{0.8}{\frac{x^{3} + sinx}{x}dx,}\)使其误差不超过\(0.1\text{x}10^{- 4}.\)

\end{document}
